\section{Stable numerical implementation of posterior summaries}
\label{app:numerical-summaries}
Family-specific integrated likelihoods, conjugate posterior updates, and complete proofs are provided directly under Propositions~\ref{prop:gaussian-family}--\ref{prop:beta-response-family}; the appendix does not repeat those calculations. Here we record two implementation details whose incorrect handling changes posterior inference even when every family-specific segment calculation is correct.

\subsection{A separate partition-prior normalizer}
For a fixed segment count, the unnormalized prior score of a path is $\prod_q\gamma_u(t_{q-1},t_q)$. Defining $C_0(0)=1$ and $C_0(j>0)=0$, its normalizer obeys
\begin{equation}
C_q(j)=\sum_{i=q-1}^{j-1}C_{q-1}(i)\gamma_u(i,j),
\qquad C_k(u)=C_k(n).
\label{eq:prior-normalizer-recurrence}
\end{equation}
Under $\gamma_u\equiv1$ with all ordered paths allowed, $C_k(u)=\binom{n-1}{k-1}$. That normalizer is never one merely because the local factors equal one, except in the degenerate endpoint counts. Excluded paths receive $\gamma_u(i,j)=0$; likelihood-zero paths are not removed from the \emph{prior} sum.

\begin{algorithm}[H]
\caption{Log-space normalization of a factorized partition prior}
\label{alg:partition-prior-normalizer}
\KwIn{Log prior segment factors $g_{ij}=\log\gamma_u(i,j)$ with $-\infty$ for structurally excluded blocks, and $k_{\max}$.}
Initialize $P_{0,0}=0$ and $P_{0,j>0}=-\infty$\;
\For{$q=1$ \KwTo $k_{\max}$}{
  \For{$j=q$ \KwTo $n$}{
    $P_{q,j}\leftarrow\logsumexp_{i=q-1}^{j-1}(P_{q-1,i}+g_{ij})$\;
  }
}
\KwOut{$\log C_k(u)=P_{k,n}$ for each feasible $k$; reject $k$ if $C_k=0$.}
\end{algorithm}

\subsection{Signed first moments and block inclusion weights}
Let $m^{(r)}_{ij}$ denote the posterior block moment $\mathbb E[m_\star(\theta)^r\mid y_{i+1:j}]$. In the forward--backward decomposition, the posterior probability that $(i,j]$ is the segment covering a specified coordinate under fixed $k$ is
\begin{equation}
\omega_{ij}(k)=\frac{\widetilde A_{ij}^{(0)}}{L_{k,n}}
\sum_{q=1}^{k}L_{q-1,i}R_{k-q,j}.
\label{eq:block-inclusion-weight}
\end{equation}
For coordinate $s$ the first and second moments are $\sum_{i<s\le j}\omega_{ij}(k)m^{(r)}_{ij}$, $r=1,2$. A Gaussian posterior mean can be negative, so $\log A^{(1)}_{ij}$ is not generally defined. The numerically stable interface stores $\log A^{(0)}_{ij}$ and signed $m^{(1)}_{ij}$ separately.

\begin{algorithm}[H]
\caption{Fixed-count Bayes curve with potentially negative segment means}
\label{alg:signed-moment-bayes-curve}
\KwIn{Forward/backward log messages, $\log\widetilde A^{(0)}_{ij}$, posterior segment moments $m^{(1)}_{ij},m^{(2)}_{ij}$, fixed count $k$.}
\For{each valid segment $(i,j]$}{
  Compute $\log\omega_{ij}(k)$ with $\logsumexp_q$ using Equation~\eqref{eq:block-inclusion-weight}\;
  Retain the sign of $m^{(1)}_{ij}$ separately from $\log|m^{(1)}_{ij}|$\;
}
\For{each coordinate $s$}{
  Sum positive and negative contributions $\omega_{ij}(k)m^{(1)}_{ij}$ for $i<s\le j$ using compensated arithmetic or signed-log-sum-exp\;
  Sum nonnegative contributions $\omega_{ij}(k)m^{(2)}_{ij}$ with stabilized arithmetic\;
  Set $\widehat\mu_s$ to the signed first-moment sum and $v_s$ to the second-moment sum minus $\widehat\mu_s^2$\;
  Check $v_s\ge-\delta$ for declared numerical tolerance $\delta$, and report any failure\;
}
\KwOut{Fixed-$k$ Bayes posterior-mean curve and pointwise posterior variances.}
\end{algorithm}

\paragraph{Masking rule.}
An unobserved outcome is absent from every likelihood and sufficient statistic. In a shared-boundary analysis, a subject with no observed outcomes in one block contributes log segment evidence zero when it has a proper segment prior. It must not contribute $-\infty$ merely because that subject supplies no information. Minimum-observation requirements are structural admissibility constraints and are represented in the prior support instead.

\section{Log-space dynamic programs}
\label{app:algorithms}

\begin{algorithm}[H]
\caption{Log-space sum-product messages}
\label{alg:log-sum-product-app}
\KwIn{Prior-adjusted log scores $\widetilde\ell_{ij}$; $k_{\max}$.}
Initialize $\Lambda_{0,0}=0$, $\Lambda_{0,j>0}=-\infty$, $\mathcal R_{0,n}=0$, and $\mathcal R_{0,i<n}=-\infty$\;
\For{$k=0$ \KwTo $k_{\max}-1$}{
  \For{$j=k+1$ \KwTo $n$}{
    $\Lambda_{k+1,j}\leftarrow\logsumexp_{h=k}^{j-1}\{\Lambda_{k,h}+\widetilde\ell_{h,j}\}$\;
  }
  \For{$i\leftarrow n-k-1,n-k-2,\ldots,0$}{
    $\mathcal R_{k+1,i}\leftarrow\logsumexp_{h=i+1}^{n-k}\{\widetilde\ell_{i,h}+\mathcal R_{k,h}\}$\;
  }
}
Normalize $\log p(k)+\Lambda_{k,n}-\log C_k(u)$ over $k$\;
\KwOut{$\Lambda$, $\mathcal R$, and $p(k\mid y)$.}
\end{algorithm}

\begin{algorithm}[H]
\caption{Joint-MAP max-sum and backtracking}
\label{alg:max-sum-app}
\KwIn{Prior-adjusted log scores $\widetilde\ell_{ij}$; fixed $k$.}
Set $M_{0,0}=0$ and $M_{0,j>0}=-\infty$\;
\For{$q=1$ \KwTo $k$}{
  \For{$j=q$ \KwTo $n$}{
    $h^\star(q,j)\leftarrow\arg\max_{h=q-1}^{j-1}\{M_{q-1,h}+\widetilde\ell_{h,j}\}$ under the declared tie rule\;
    $M_{q,j}\leftarrow M_{q-1,h^\star(q,j)}+\widetilde\ell_{h^\star(q,j),j}$\;
  }
}
Set $\widehat t_k=n$ and recursively $\widehat t_{q-1}=h^\star(q,\widehat t_q)$\;
\KwOut{$\widehat t$ and terminal score $M_{k,n}$.}
\end{algorithm}

The recommended numerical checks are
\begin{align*}
\sum_{k=1}^{k_{\max}}p(k\mid y)&=1,\\
|\Lambda_{k,n}-\mathcal R_{k,0}|&\le\delta,\\
\sum_{h=1}^{n-1}p(b_h=1\mid y,k)&=k-1,\\
\left|M_{k,n}-\sum_{q=1}^{k}
\widetilde\ell_{\widehat t_{q-1},\widehat t_q}\right|&\le\delta.
\end{align*}
with a documented floating-point tolerance $\delta$.

\section{Supplementary agreement diagnostics}
\label{app:agreement}

The tables below preserve executed comparator values while making the reference semantics explicit. They are not external-accuracy tables.

\begin{table}[H]
\centering
\small
\caption{Well-log boundary agreement with the \BayesBreak MAP boundary set at tolerance three. No independently verified geological boundary annotation was used.}
\label{tab:welllog-agreement-app}
\begin{tabular}{@{}lrr@{}}
\toprule
Method & Agreement $F_1@3$ & Interpretation \\
\midrule
Fearnhead-configured reference DP & 0.93 & agreement with BB MAP \\
Optimal partitioning & 0.91 & agreement with BB MAP \\
Binary segmentation & 0.86 & agreement with BB MAP \\
PELT & 0.18 & agreement with BB MAP \\
Wild binary segmentation & 0.05 & agreement with BB MAP \\
\bottomrule
\end{tabular}
\end{table}

\begin{table}[H]
\centering
\small
\caption{S\&P 500 and methylation comparator agreement. Segment-count restrictions were matched for the indicated comparisons.}
\label{tab:spx-methyl-agreement-app}
\begin{tabularx}{0.96\textwidth}{@{}L{0.20\textwidth}L{0.28\textwidth}R{0.11\textwidth}R{0.13\textwidth}Y@{}}
\toprule
Dataset & Method & $\widehat k$ & Agreement $F_1@3$ & Scope \\
\midrule
S\&P 500 & optimal partitioning, matched $k$ & 29 & 0.82 & vs BB Gaussian MAP \\
S\&P 500 & PELT & 79 & 0.53 & vs BB Gaussian MAP \\
S\&P 500 & binary segmentation, matched $k$ & 29 & 0.43 & vs BB Gaussian MAP \\
S\&P 500 & wild binary segmentation, matched $k$ & 29 & 0.07 & vs BB Gaussian MAP \\
S\&P 500 & Fearnhead-configured reference DP & 13 & 0.55 & vs BB Gaussian MAP \\
Methylation & optimal partitioning, matched $k$ & 15 & 0.79 & vs BB Beta-valued MAP \\
Methylation & binary segmentation, matched $k$ & 15 & 0.57 & vs BB Beta-valued MAP \\
Methylation & PELT & 5 & 0.44 & vs BB Beta-valued MAP \\
Methylation & wild binary segmentation, matched $k$ & 15 & 0.14 & vs BB Beta-valued MAP \\
\bottomrule
\end{tabularx}
\end{table}

The incompatible array-CGH comparator output is intentionally absent from these tables. Its flattened 95,245-value axis and the 2,215-probe boundary axis do not define a common matching problem.

\section{Reproducibility}
\label{app:reproducibility}
Reproduction of an experiment requires specifying the exact data source and
coordinate system, the preprocessing and exclusion rules, the observation
model and its descriptors, the prior, the numerical method, the random seeds,
and the model-selection and evaluation criteria. The prospective simulations
and their corresponding figures and tables are specified in
Section~\ref{sec:experimental-design}; full machine-readable configurations
and run instructions belong with the accompanying code.
